% GENERATED FILE. Edit canonical lessons, then run npm run generate:books.
% canonical-source-hash: fnv1a64:af771318822e692f
% canonical-lessons: UR-C67-registan, UR-C67-kutta, UR-C67-billi, UR-C67-ghora, UR-C67-gae

\chapter{Desert, Dog, Cat, Horse, Cow}
\label{ch:ur-desert-dog-cat-horse-cow}
\hlchaptermodality{67}

\begin{chapteropening}
\textbf{By the end of this chapter:} \emph{I can say a desert, a dog, a cat, a horse and a cow.}

Complete the last lesson of chapter 67: I can say a desert, a dog, a cat, a horse and a cow.
\end{chapteropening}

\section[\texorpdfstring{\ur{ریگستان} (registān)}{ریگستان (registān)}]{\ur{ریگستان} (registān) — a desert}

\label{lesson:UR-C67-registan}



\begin{quote}
Before the new one: say the Urdu for fire, then the Urdu for snow; ice.
\end{quote}

\subsection*{You'll want to know: \ur{ریگستان}}
\textbf{\ur{ریگستان}} — \emph{registān} — \textquotedblleft{}a desert\textquotedblright{}.

\begin{cousinweb}
\textbf{\ur{ریگستان}} (\emph{registān}) is a loan, usually traced to Persian \emph{rēgistān}, \textquotedblleft{}desert\textquotedblright{}. It is built from \emph{rēg}, \textquotedblleft{}sand\textquotedblright{}, and \emph{-stān}, \textquotedblleft{}place\textquotedblright{}.
\end{cousinweb}

\begin{grammarlens}[title={What you've built}]
One more word for this chapter.
\end{grammarlens}

\subsection*{Your turn}
\begin{itemize}
\raggedright
  \item \emph{Say it:} \emph{registān}
  \item \emph{Say it:} \emph{registān}, once more
  \item \emph{Say it:} say \emph{registān}
  \item \emph{Recall:} say the Urdu for fire, then the Urdu for snow; ice, then say \emph{registān} again
\end{itemize}

\begin{tcolorbox}[breakable,colback=teal!4,colframe=teal!35!black,title={Before you move on}]
What does \ur{ریگستان} mean? (\textquotedblleft{}A desert\textquotedblright{}.) Say it once more.
\end{tcolorbox}
\section[\texorpdfstring{\ur{کتا} (kuttā)}{کتا (kuttā)}]{\ur{کتا} (kuttā) — a dog}

\label{lesson:UR-C67-kutta}



\begin{quote}
Before the new one: say the Urdu for snow; ice, then the Urdu for a desert.
\end{quote}

\subsection*{You'll want to know: \ur{کتا}}
\textbf{\ur{کتا}} — \emph{kuttā} — \textquotedblleft{}a dog\textquotedblright{}.

\begin{cousinweb}
\textbf{\ur{کتا}} (\emph{kuttā}) is a tadbhava, a word inherited and reshaped by sound change. It comes from \emph{*kutta}, \textquotedblleft{}dog\textquotedblright{}, a form CDIAL reconstructs from its descendants (CDIAL 3275). CDIAL calls it imitative, a word made from the sound used to call a dog.
\end{cousinweb}

\begin{grammarlens}[title={What you've built}]
One more word for this chapter.
\end{grammarlens}

\subsection*{Your turn}
\begin{itemize}
\raggedright
  \item \emph{Say it:} \emph{kuttā}
  \item \emph{Say it:} \emph{kuttā}, once more
  \item \emph{Say it:} say \emph{kuttā}
  \item \emph{Recall:} say the Urdu for snow; ice, then the Urdu for a desert, then say \emph{kuttā} again
\end{itemize}

\begin{tcolorbox}[breakable,colback=teal!4,colframe=teal!35!black,title={Before you move on}]
What does \ur{کتا} mean? (\textquotedblleft{}A dog\textquotedblright{}.) Say it once more.
\end{tcolorbox}
\section[\texorpdfstring{\ur{بلی} (billī)}{بلی (billī)}]{\ur{بلی} (billī) — a cat}

\label{lesson:UR-C67-billi}



\begin{quote}
Before the new one: say the Urdu for a desert, then the Urdu for a dog.
\end{quote}

\subsection*{You'll want to know: \ur{بلی}}
\textbf{\ur{بلی}} — \emph{billī} — \textquotedblleft{}a cat\textquotedblright{}.

\begin{cousinweb}
\textbf{\ur{بلی}} (\emph{billī}) is inherited. CDIAL lists Hindi \emph{billī}, glossed \textquotedblleft{}kitten\textquotedblright{}, under a reconstructed \emph{*billa}, a form it sets beside Sanskrit \emph{biḍāla}, \textquotedblleft{}cat\textquotedblright{} (CDIAL 9237).
\end{cousinweb}

\begin{grammarlens}[title={What you've built}]
One more word for this chapter.
\end{grammarlens}

\subsection*{Your turn}
\begin{itemize}
\raggedright
  \item \emph{Say it:} \emph{billī}
  \item \emph{Say it:} \emph{billī}, once more
  \item \emph{Say it:} say \emph{billī}
  \item \emph{Recall:} say the Urdu for a desert, then the Urdu for a dog, then say \emph{billī} again
\end{itemize}

\begin{tcolorbox}[breakable,colback=teal!4,colframe=teal!35!black,title={Before you move on}]
What does \ur{بلی} mean? (\textquotedblleft{}A cat\textquotedblright{}.) Say it once more.
\end{tcolorbox}
\section[\texorpdfstring{\ur{گھوڑا} (ghoṛā)}{گھوڑا (ghoṛā)}]{\ur{گھوڑا} (ghoṛā) — a horse}

\label{lesson:UR-C67-ghora}



\begin{quote}
Before the new one: say the Urdu for a dog, then the Urdu for a cat.
\end{quote}

\subsection*{You'll want to know: \ur{گھوڑا}}
\textbf{\ur{گھوڑا}} — \emph{ghoṛā} — \textquotedblleft{}a horse\textquotedblright{}.

\begin{cousinweb}
\textbf{\ur{گھوڑا}} (\emph{ghoṛā}) is a tadbhava, a word inherited from Sanskrit and reshaped by sound change. It comes from \emph{ghoṭa}, \textquotedblleft{}horse\textquotedblright{} (CDIAL 4516). CDIAL calls the word non-Aryan, probably Dravidian.
\end{cousinweb}

\begin{grammarlens}[title={What you've built}]
One more word for this chapter.
\end{grammarlens}

\subsection*{Your turn}
\begin{itemize}
\raggedright
  \item \emph{Say it:} \emph{ghoṛā}
  \item \emph{Say it:} \emph{ghoṛā}, once more
  \item \emph{Say it:} say \emph{ghoṛā}
  \item \emph{Recall:} say the Urdu for a dog, then the Urdu for a cat, then say \emph{ghoṛā} again
\end{itemize}

\begin{tcolorbox}[breakable,colback=teal!4,colframe=teal!35!black,title={Before you move on}]
What does \ur{گھوڑا} mean? (\textquotedblleft{}A horse\textquotedblright{}.) Say it once more.
\end{tcolorbox}
\section[\texorpdfstring{\ur{گائے} (gāe)}{گائے (gāe)}]{\ur{گائے} (gāe) — a cow}

\label{lesson:UR-C67-gae}



\begin{quote}
Before the new one: say the Urdu for a cat, then the Urdu for a horse.
\end{quote}

\subsection*{You'll want to know: \ur{گائے}}
\textbf{\ur{گائے}} — \emph{gāe} — \textquotedblleft{}a cow\textquotedblright{}.

\begin{cousinweb}
\textbf{\ur{گائے}} (\emph{gāe}) is inherited from Sanskrit \emph{gāvī}, \textquotedblleft{}cow\textquotedblright{}, a form of \emph{go}, \textquotedblleft{}cow, ox\textquotedblright{} (CDIAL 4147).
\end{cousinweb}

\begin{grammarlens}[title={What you've built}]
One more word for this chapter.
\end{grammarlens}

\subsection*{Your turn}
\begin{itemize}
\raggedright
  \item \emph{Say it:} \emph{gāe}
  \item \emph{Say it:} \emph{gāe}, once more
  \item \emph{Say it:} say \emph{gāe}
  \item \emph{Recall:} say the Urdu for a cat, then the Urdu for a horse, then say \emph{gāe} again
\end{itemize}

\begin{tcolorbox}[breakable,colback=teal!4,colframe=teal!35!black,title={Before you move on}]
What does \ur{گائے} mean? (\textquotedblleft{}A cow\textquotedblright{}.) Say it once more.
\end{tcolorbox}
